\section{Metodología experimental}

\subsection{Plataformas}

El entorno local de desarrollo es WSL2 con Ubuntu 24.04 x86-64 sobre un AMD Ryzen 5 3600 de seis núcleos y doce hilos, 7,7~GiB de memoria y 16~MiB de caché L3. Se utilizó GCC 13.3.0 y Open MPI 4.1.6. El proyecto se construye principalmente con CMake y mantiene un Makefile portable para el entorno académico.

El programa se escribió en C11 y se compiló con advertencias estrictas. Durante el desarrollo también se revisaron el formato, posibles errores de memoria y usos incorrectos de variables. Los comandos principales de compilación y ejecución se conservan en el archivo \texttt{README.md}.

Las mediciones de FING se realizaron el 9 de agosto de 2026 en los equipos
\texttt{pcunix40} y \texttt{pcunix42}. Cada equipo posee un Intel Core i3-4170
a 3,70~GHz, dos núcleos físicos, cuatro hilos lógicos, 15~GiB de memoria y
3~MiB de caché L3. Ambos ejecutaban Fedora Linux con kernel 6.19.14. Se utilizó
GCC 15.2.1 y MPICH 4.2.2, habilitado mediante el módulo
\texttt{mpi/mpich-x86\_64}. Los trabajos se lanzaron directamente por SSH sobre
los dos equipos disponibles, sin un planificador. El equipo \texttt{lulu} se
utilizó únicamente como punto de acceso. MPICH se ejecutó con el proveedor TCP
para utilizar la red Ethernet de las pcunix.

\subsection{Configuraciones y procedimiento}

Para medir escalabilidad fuerte se mantiene fija la instancia y se cambia la cantidad de procesos MPI e hilos OpenMP. Los checkpoints se desactivan y todas las configuraciones deben terminar con el mismo hash.

Para observar el cambio espacial se reconstruye la grilla inicial con la misma semilla y se lee el estado final guardado por la simulación. Cada celda se pinta según su clase social, o como vivienda vacía o celda no residencial. Además se calcula, para cada hogar no aislado, la proporción ponderada de vecinos de su misma clase. El promedio de esas proporciones permite comparar el agrupamiento inicial y final con un único valor. Se usan el mismo radio y los mismos pesos del modelo.

La validación funcional usa tres tamaños: mínimo de $8\times6$, mediano de $256\times160$ y funcional de $1024\times640$. Las primeras mediciones de rendimiento también usaron el escenario funcional. En ellas se ejecutó un calentamiento y cinco repeticiones por configuración, y se informó la mediana.

La instancia principal \texttt{hpc.conf} tiene $4096\times2560$ celdas y 240 iteraciones. Mantiene la distribución social, las preferencias y la regla económica; usa 20\% de viviendas vacías, permanencia cero, financiación de 60\%, plazo de 360 meses y un límite de cuota de 40\% del ingreso. La grilla es una carga sintética y no una estimación de la cantidad de viviendas de Montevideo.

Debido a la duración de la instancia grande, la campaña principal realiza una ejecución de cada configuración. Las mediciones más pequeñas y repetidas se conservan como referencia de variación. Se comparan la versión secuencial y las configuraciones $1\times1$, $1\times2$, $1\times4$, $2\times1$, $2\times2$ y $2\times4$. El primer número indica procesos MPI y el segundo, hilos OpenMP por proceso.

\subsection{Métricas}

Para $T_1$ secuencial y $T_p$ paralelo, se calculan
\begin{equation}
S_p=\frac{T_1}{T_p}, \qquad E_p=\frac{S_p}{p}.
\end{equation}
También se miden tiempos máximos por fase, desbalance $L=\max_i W_i/\overline{W}$, volumen lógico de comunicación y métricas funcionales del modelo. El trabajo $W_i$ es la cantidad de hogares elegibles cuya búsqueda de destino pertenece al proceso $i$.

\subsection{Reproducibilidad}

El script \texttt{scripts/ejecutar\_experimentos.py} crea un directorio por ejecución, inicia las versiones secuencial e híbrida, compara los hashes y produce \texttt{resumen.csv}. Fija \texttt{OMP\_PLACES=cores}, \texttt{OMP\_PROC\_BIND=close} y la cantidad de hilos de cada configuración. Los archivos de la campaña principal se guardan en \texttt{informe/datos/hpc\_final} junto con la configuración necesaria para repetirla.

El proyecto se compila con CMake y cada versión recibe un archivo mediante \texttt{--config}. La ejecución secuencial usa \texttt{schelling\_seq}; la híbrida se inicia con \texttt{mpirun}, indicando procesos e hilos. La salida incluye la configuración efectiva, métricas y tiempos por iteración, datos de comunicación y el estado final. Con \texttt{hpc.conf} y semilla 42, todas las configuraciones finalizaron con el hash \texttt{6b39928d840ed346}.
